\section{Crossing information and spatial growth}\label{cross:sec:crossing}
Remote feedback is a temporal requirement on persistent nontranslation. Crossing
records give an additional spatial requirement. They describe the information
carried by the head through a cut, a method introduced for Turing machines by
Hennie [7]. The following argument proves the needed half-plane coupling directly.
Its local incoming-axis invariant is standard ant geometry [2].
The log-log work-space threshold reviewed in [8] concerns input length; here
the independent variable is the elapsed time of one complete ant orbit.

Write $\ETC$ for eventual translation of position, heading and turn with a
nonzero drift and a fixed period. The complete black field may retain a growing
wake. This is the same translating alternative used in the feedback theorems.

For real $\alpha,\beta\ne0$, normalize
$\min(|\alpha|,|\beta|)=1$ and define
\[
 w_{\alpha,\beta}(T)=
 \max_{0\le t\le T}(\alpha p_t.x+\beta p_t.y)
 -\min_{0\le t\le T}(\alpha p_t.x+\beta p_t.y).
\]
Here $p_t.x$ and $p_t.y$ denote the horizontal and vertical coordinates of $p_t$.
This is normal-coordinate width, not Euclidean width when the normal has a
large norm. Both components must be nonzero.

\begin{theorem}[Width gap]\label{cross:thm:gap}
Every classical finite-support orbit that is not $\ETC$ satisfies, for every
such normal,
\[
 \liminf_{T\to\infty}\frac{w_{\alpha,\beta}(T)}{\log\log T}
 \ge\frac1{\log2}.
\]
The finite bound from which this follows is uniform in the normalized normal.
Consequently the conclusion also holds for time-dependent normal choices,
and for the infimum of these widths over all admissible normals at each time.
\end{theorem}
\subsection{A finite certificate with its own future protection}
For an integer $a$ write $R_a=\{(x,y):x>a\}$. Suppose two times $s<t$ and cuts
$a<b$ satisfy
\begin{align}
 p_t&=p_s+d,\quad h_t=h_s,\quad d=(b-a,\delta),\label{cross:eq:poses}\\
 (B_s\cap R_a)+d&=B_t\cap R_b,\label{cross:eq:field}\\
 p_u&\in R_a\quad(s\le u\le t).\label{cross:eq:guard}
\end{align}
The comparison in \eqref{cross:eq:field} is the complete inside field.

\begin{lemma}[Nested-domain pumping]\label{cross:lem:nested}
Conditions \eqref{cross:eq:poses}--\eqref{cross:eq:guard} imply that the trace from $s$ is
permanently translating with period $t-s$ and drift $d$.
\end{lemma}
\begin{proof}
Until its first exit, the trajectory in a half-plane depends only on its inside
field and starting pose. Thus the two inside states have the same first-exit
duration $\tau\in\N\cup\{\infty\}$. Put $P=t-s$. If $\tau\le P$, the finite
guard fails. If $P<\tau<\infty$, the orbit must exit the smaller $R_b$ between
$t$ and its exit from $R_a$ at $s+\tau$. But the second inside state predicts
its first exit at $t+\tau>s+\tau$. Therefore $\tau=\infty$. Determinism now
identifies the entire two futures up to translation. Since $b>a$, the drift
is nonzero. No exterior cell has been removed from the physical state.
\end{proof}

Set
\[
 a_0=\max\bigl(\{p_0.x\}\cup\{q.x:q\in B_0\}\bigr).
\]
For a finite actual prefix through $T$, the crossing record of the cut
$x=c+\tfrac12$ lists every crossing in chronological order, with its direction
and height. Normalize heights by subtracting their minimum. If $c\ge a_0$ and
$p_T.x>c$, the directions are necessarily $E,W,\ldots,E$.

\begin{lemma}[Complete crossing records]\label{cross:lem:records}
Two equal nonempty normalized records at cuts $a_0\le a<b<p_T.x$ produce a
certificate satisfying Lemma~\ref{cross:lem:nested} at times no later than $T$.
\end{lemma}
\begin{proof}
Both right half-planes are initially white. Their first complete right-side
excursions have translated entry poses and therefore translated read/write
histories until the first exit. Between excursions, left-side updates do not
change the right-side field. Induction over the equal records gives translated
complete inside fields at their last eastward entries, at times $s$ and $t$.
The translation in height is the difference of record minima. The last entry
at $b$ must follow the last entry at $a$, since otherwise ending beyond $b$
would require a later $b$-entry. There is no exit from $R_a$ between $s$ and
$t$, because $s$ is its last entry before $T$. All three conditions follow.
\end{proof}
The records are finite histories, not a query to the infinite future. A far
initial black cell to the right invalidates the hypothesis even when it has
not yet altered the observed finite path.

\subsection{Exact record capacity and crossing work}
Number headings $N,E,S,W$ by $0,1,2,3$. The value $h-x-y\pmod2$ is invariant.
An east crossing of cut $c+\tfrac12$ arrives in column $c+1$; a west crossing
arrives in column $c$. Both arrivals have horizontal heading, so their height
parities differ. Consecutive heights in a crossing record therefore alternate
parity, also after normalization.

For $H\ge0$ put
\[
 E=\lfloor H/2\rfloor+1,\qquad O=\lceil H/2\rceil,
 \qquad U=EO,\qquad V=(E-1)O.
\]
\begin{proposition}[Parity capacity]\label{cross:prop:count}
The number of normalized height words of length $2j+1$, with span at most $H$
and alternating parity, is
\[
 c_H(j)=(H+1)U^j-HV^j,\qquad j\ge0,
\]
where $0^0=1$.
\end{proposition}
\begin{proof}
Over $\{0,\ldots,H\}$ the count before requiring minimum zero is
$E^{j+1}O^j+O^{j+1}E^j=(H+1)U^j$. Words omitting zero use $E-1$ even symbols
and $O$ odd symbols, giving $HV^j$. Subtract them.
\end{proof}
Not every word counted by this proposition is necessarily physically
realizable. It is an exact count of a necessary combinatorial class.

\paragraph{Incoming requests reconstruct the record.}
For an eligible cut the genuine record has the form
\[
 E(y_0),W(z_0),E(y_1),W(z_1),\ldots,E(y_j).
\]
The exit heights $z_i$ are determined responses to the incoming heights $y_i$.
This causal constraint gives a smaller capacity than the parity count.

\begin{lemma}[Request injectivity]\label{cross:lem:requests}
The incoming heights $(y_0,\ldots,y_j)$ determine every completed right-side
excursion: its exit height, duration and complete inside field after return.
They also determine the complete inside field immediately before the final
entry. Vertically translated request lists give translated full records.
\end{lemma}
\begin{proof}
The first entry sees an entirely white half-plane. Literal dynamics uniquely
determine its first exit and all writes up to that exit, when it exists.
During an outside excursion no inside cell changes. The next entry therefore
sees exactly the retained field produced by the preceding calls. Induction
proves the assertion for all completed excursions. Translation invariance
also identifies histories at horizontally shifted cuts. The final request
records its entry; it need not return.
\end{proof}
All entry heights have the same parity. Put $m=\min_i y_i$ and encode the
record by the integer word
\[
 \rho_c=((y_0-m)/2,\ldots,(y_j-m)/2).
\]
Two genuine records agree up to translation if and only if these request codes
agree: restriction to the entries gives one direction, and
Lemma~\ref{cross:lem:requests} gives the other. Thus equal request codes at two
eligible cuts yield the certificate of Lemma~\ref{cross:lem:records}.

\begin{proposition}[Request capacity]\label{cross:prop:requests}
For $H\ge1$, put $q=\lfloor H/2\rfloor+1$. The number of genuine normalized
records of length $2j+1$ and span at most $H$ is at most
\begin{equation}
 b_H(j)=q^{j+1}-(q-1)^{j+1}.\label{cross:eq:request-capacity}
\end{equation}
For $H=0$, only the one-letter signature exists.
\end{proposition}
\begin{proof}
A request code has $j+1$ letters in $\{0,\ldots,q-1\}$ and contains zero.
Subtracting the words omitting zero gives \eqref{cross:eq:request-capacity}.
Request injectivity makes this an upper bound for genuine signatures. At
height zero a completed return would violate the opposite parities of entry
and exit heights.
\end{proof}
For example, at $H=10$ and $j=5$ the parity capacity is $169{,}643{,}750$,
whereas the request capacity is $31{,}031$. These count combinatorial classes;
a request code contributes an available object even when it has no returning,
span-admissible physical response.

Write $A_H^*(r)=\sum_{j=0}^r b_H(j)$, with $A_H^*(-1)=0$ and
$b_0(j)=\mathbf1_{j=0}$. Let $\mathcal P_H^*(N)$ be the minimum total length
of $N$ distinct objects, with $b_H(j)$ objects of length $2j+1$.
If $r$ is the first index with $A_H^*(r)\ge N$, shortest-layer filling gives
\begin{equation}
 \mathcal P_H^*(N)=\sum_{j<r}(2j+1)b_H(j)
 +(2r+1)\bigl(N-A_H^*(r-1)\bigr).\label{cross:eq:packing}
\end{equation}
Set $\mathcal P_H^*(0)=0$ and $\mathcal P_0^*(N)=\infty$ for $N>1$.
At height one, $b_1(j)=1$, so $\mathcal P_1^*(N)=N^2$.
For $q\ge2$, the cumulative capacity is the integer geometric sum
\[
 A_H^*(r)=\frac{q(q^{r+1}-1)}{q-1}
 -(q-1)\sum_{j=0}^r(q-1)^j.
\]

\begin{theorem}[Request-work dichotomy]\label{cross:thm:work}
Let $N=(p_T.x-a_0)_+$. For $N>0$ let $H$ be the maximum span of the crossing
records at the $N$ cuts $a_0,\ldots,p_T.x-1$. Either a permanent translation
certificate occurs in this finite prefix or
\[
 \mathcal P_H^*(N)\le\lceil T/2\rceil.
\]
\end{theorem}
\begin{proof}
Without a certificate, the $N$ request codes are distinct. Request capacity
therefore bounds their summed full-record lengths below by
\eqref{cross:eq:packing}. Each listed crossing is a different horizontal update,
and horizontal and vertical updates alternate.
\end{proof}

\begin{corollary}[Entropy form]\label{cross:cor:entropy}
In the second alternative, when $H\ge2$ and $N>0$, let
$\overline J=\lfloor(\lceil T/2\rceil-N)/2\rfloor$. Then
\begin{equation}
 \log(N!)\le\overline J\log q
       +N\log\bigl(2(1+\overline J/N)\bigr).
 \label{cross:eq:entropy}
\end{equation}
\end{corollary}
\begin{proof}
The difference-of-powers identity gives $b_H(j)\le(j+1)q^j$, and hence
$A_H^*(r)\le2(r+1)q^r$ for $q\ge2$. Order the pair counts of the $N$ actual
records as $j_1\le\cdots\le j_N$ and put $J=\sum_i j_i$. Since
$i\le A_H^*(j_i)$, summing logarithms and applying concavity gives
\[
 \log(N!)\le J\log q+N\log\bigl(2(1+J/N)\bigr).
\]
Their total length is $2J+N\le\lceil T/2\rceil$, so $J\le\overline J$.
The right-hand side increases with $J$.
\end{proof}

\begin{corollary}[Speed--spread tradeoff]\label{cross:cor:speed}
For an orbit which is not $\ETC$, suppose $T_i\to\infty$ and
$N(T_i)/T_i\to v>0$. Then $v<1/2$ and its record spans satisfy
\[
 \liminf_i\frac{\log(H(T_i)+2)}{\log T_i}\ge\frac{4v}{1-2v}.
\]
\end{corollary}
\begin{proof}
The cases $H\le1$ give at least quadratic work or exclude $N>1$, and are
therefore eventually absent. Also $H\le\lceil T/2\rceil$, since height
variation cannot exceed the number of vertical moves. The work inequality
already gives $v\le1/2$. If $v=1/2$, then $\overline J=o(T)$, and the
right-hand side of \eqref{cross:eq:entropy} is $o(T\log T)$, contradicting
$\log(N!)\sim vT\log T$. For $v<1/2$, divide by $T\log T$ and use
$\overline J/T\to(1-2v)/4$ and $\log N/\log T\to1$.
\end{proof}
Half-turns and quarter-turns give the other directions. For example, speed
$v=1/52$ requires exponent at least $2/25$ on a nontranslating orbit.
Subpolynomial crossing-height spread is consequently incompatible with
positive horizontal upper speed. Together with $H=O(T)$ the bound also gives
$v\le1/6$. No bound on repeated visits is assumed.

\subsection{A numerical finite-profile deadline}
For $H\ge0$, $C\ge1$, and a complete source, define
\begin{align*}
 R&=\max\bigl(\{0\}\cup\{|q.x-p_0.x|:q\in B_0\}\bigr),\\
 L&=H+1,\qquad M(H,C)=A_H^*(LC-1),\\
 D(R,H,C)&=\bigl[2(R+M(H,C))+1\bigr]LC.
\end{align*}
\begin{theorem}[Profile deadline]\label{cross:thm:profile}
Suppose a finite prefix through $T$ has vertical span at most $H$ in every
visited column and at most $C$ updates of every site. Either it contains a
permanent translation certificate or $T\le D(R,H,C)$.
\end{theorem}
\begin{proof}
In an eligible record, east crossings update their departure sites in its left
column. There are at most $LC$ such updates. Since the record starts and ends
east, its odd length is at most $2LC-1$. Every such record therefore lies in a
class of size $M(H,C)$. Reaching a point with
$x>p_0.x+R+M(H,C)$ would give more than $M(H,C)$ cuts beyond the complete initial
support at that subprefix endpoint, forcing a duplicate record. A half-turn
gives the analogous left bound. If no certificate occurs, every head position
has $|x-p_0.x|\le R+M(H,C)$. There are at most $2(R+M)+1$ visited columns,
each with at most $L$ updated sites and at most $C$ updates per site.
\end{proof}
The final unupdated head is not charged as an update. At a final east crossing,
its departure is still in the counted left column; all west crossings of an
eligible record are earlier. Black cells outside the corridor remain present.
Their coordinates matter through the initial frontier, not through an assumed
future-active quotient. The bound is linear in numerical $R$ at fixed $H,C$,
but is not a polynomial bit-complexity claim.

\subsection{Local visit bounds and slowly widening trajectories}
Let $n(v)$ count updates in a finite segment, excluding the terminal position.
For either adjacent pair $u=v+(\pm1,0)$ and $z=v+(0,\pm1)$,
\begin{equation}
 n(v)\le2n(u)+2n(z)+5.\label{cross:eq:corner}
\end{equation}
For completeness, the incoming axis at $v$ is fixed. A visit is exceptional if
it is the initial update, arrives from the selected incoming neighbour, or
leaves to the selected outgoing neighbour. On every other visit both directions
are fixed, hence the turn and read color are fixed. Alternation of colors forces
at least $\lfloor n(v)/2\rfloor$ exceptions. There are at most $n(u)+n(z)+2$,
including the initial and terminal allowances, proving \eqref{cross:eq:corner}.

If all head positions of a prefix lie in $a\le\alpha x+\beta y\le b$ with
$\min(|\alpha|,|\beta|)=1$, choose the adjacent pair increasing the linear form.
Both children increase it by at least one. The analogous pair decreases it.
Iterating toward the nearer boundary, with $F(-1)=0$ and $F(d)=4F(d-1)+5$, gives
\[
 n(v)\le C_h:=\frac{5(4^{\lfloor h/2\rfloor+1}-1)}3,
 \qquad h=\lfloor b-a\rfloor.
\]
Each column contains at most $\lfloor(b-a)/|\beta|\rfloor+1\le h+1$ lattice
sites. Therefore Theorem~\ref{cross:thm:profile} proves the explicit alternative
\begin{equation}
 \text{a certificate occurs, or}\qquad
 T\le D_{\mathrm{obl}}(R,h):=D(R,h,C_h).\label{cross:eq:deadline}
\end{equation}
No infinite confinement is assumed in \eqref{cross:eq:deadline}.

\begin{proof}[Proof of Theorem~\ref{cross:thm:gap}]
For $h\ge2$, $C_h\le(20/3)2^h$ and $q\le h+1$. Bounding $A_h^*$ by
$2(r+1)q^r$ at $r=(h+1)C_h-1$ yields
\[
 \log D_{\mathrm{obl}}(R,h)
 =O_R\bigl((h+1)2^h\log(h+2)\bigr).
\]
On an orbit which is not $\ETC$, no finite certificate can occur. Thus for every
$T$, with $h=\lfloor w_{\alpha,\beta}(T)\rfloor$,
\[
 \log\log T\le h\log2+\log(h+1)+\log\log(h+2)+O_R(1).
\]
This first forces $h\to\infty$ and then gives the claimed liminf. All constants
are independent of the normalized normal. For an infimum choose a normal with
width less than one above the infimum; the bounded difference does not change
the asymptotic conclusion.
\end{proof}

For example, an eventual tube
\[
 |y-\theta x|\le g(|x|),\qquad \theta\ne0,
 \qquad g(r)=o\bigl(\log\log(r+e^e)\bigr)
\]
forces $\ETC$: unit speed makes its prefix normal span $o(\log\log T)$; a finite
initial part contributes only a constant. For irrational $\theta$ this is
impossible, because a nonzero integer translating drift has a nonzero linear
deviation from that slope. The width $g$ is allowed to tend to infinity.

There is a corresponding nonstraight extension. Let $f:\Z\to\Z$ be
nondecreasing with $f(x+L)\ge f(x)+1$ for some fixed $L\ge1$. A north/west path
of length $L(H+1)$ leaves the corridor $f(x)\le y\le f(x)+H$ above its top.
Expanding \eqref{cross:eq:corner} along this binary tree gives
\[
 C_{L,H}=\frac{5(4^{L(H+1)}-1)}3.
\]
The profile deadline then implies, for any nontranslating orbit whose prefix
fits that corridor with a possibly increasing $H(T)$,
\[
 \liminf_{T\to\infty}\frac{H(T)}{\log\log T}\ge\frac1{L\log4}.
\]
The bound $H(T)$ covers the entire prefix, not a window from which old deviations
have been deleted.

\subsection{Nonconstant cyclic turning rules}
The geometric information argument also applies to generalized cyclic ants.
Let a rule $\omega\in\{R,L\}^m$ contain both letters, let a read increment its
cell state modulo $m$, and let only finitely many initial states differ from
zero. Put
\[
 a=\min(\#R(\omega),\#L(\omega)),\qquad \kappa=m/a.
\]
For a cyclic prefix at phase $s$, let $n_c(s,r)$ count turn $c$. Define
\[
 b_\omega=\max_{s,c,\,0\le r<m}\{r-\kappa n_c(s,r)\}.
\]
Every full $m$-block has at least $a$ copies of either turn, so
$n\le\kappa n_c(s,n)+b_\omega$ for every prefix length. In the corner argument,
all nonexceptional visits require one fixed turn. Visits with the opposite turn
are exceptional, giving
\[
 n(v)\le\kappa(n(u)+n(z))+2\kappa+b_\omega.
\]
The rest of the proof is unchanged, replacing $C_h$ by
\[
 C_\omega(h)=\left\lceil
 (2\kappa+b_\omega)
 \frac{(2\kappa)^{\lfloor h/2\rfloor+1}-1}{2\kappa-1}
 \right\rceil.
\]

\begin{theorem}[Frequency-dependent width gap]
Every non-eventually-translating finite-support orbit for the fixed rule
$\omega$ satisfies
\[
 \liminf_{T\to\infty}\frac{w_{\alpha,\beta}(T)}{\log\log T}
 \ge\frac2{\log(2\kappa)}.
\]
\end{theorem}
\begin{proof}
Heading-axis parity and crossing parity still alternate. Complete
half-plane coupling and request reconstruction remain deterministic with $m$ colors. The local visit
bound grows as $O_\omega((2\kappa)^{h/2})$. Substitution into the same finite
profile deadline and taking two logarithms proves the claim.
\end{proof}
For classical $RL$, $\kappa=2$, $b_\omega=1$, recovering the exact earlier
local-time constant. Every balanced cyclic rule has the same asymptotic
coefficient $1/\log2$, though its finite constants may differ.

\subsection{Effectivity and reproducibility}
For a supplied rational non-axis normal and integer width bound, an implemented
algorithm follows complete states, checks the running width, and compares all
crossing records at new left/right horizontal records beyond the full initial
support. It returns a first observed width violation or a permanent translating
certificate. Without a user cap, \eqref{cross:eq:deadline} proves termination. It need
not construct the large deadline integer during execution. A capped run returns
\texttt{UNKNOWN}, never a negative dynamical conclusion.

The certificate checker independently replays the finite source, compares the
full half-plane fields, and checks the finite no-exit guard. A separate phase-ray
parity check tests all initial black hits, all self-write breakpoints, and the
final infinite interval. An unseen forward defect is therefore not ignored.
A translating answer does not assert permanent satisfaction of the supplied
width bound; a non-tangent translator may later leave it.

The Python simulator and an independently written C++ vector simulator agree
on every path position and complete endpoint field of 100 supplied inputs,
totalling 113,152 updates per implementation. Further checks exhaust the
stated small parity-word classes, verify integer packing and entropy inequalities,
check finite corner bounds, and reject malformed or unguarded certificates.
Four actual standard-tail inputs in different rotations return 104-step
certificates. These are implementation checks, not proofs that all tails are
standard. The finite cyclic-rule balance calculation is checked for every
nonconstant rule word of lengths two through eight. The general-rule classifier
is not implemented; its extension is analytic.

The request-capacity module separately checks 78 complete inputs and 86,647
updates per literal implementation. A persistent half-plane replay compares
186 complete residual fields; the controls also compare request-code and
full-signature equality, exact integer packing, and damaged histories.
It uses the same hash-pinned baseline simulator and C++ oracle as the preceding
module, together with a separate right-excursion implementation.

The reproduction entry point is
\begin{quote}\ttfamily
python3 memory\_feedback/reproduce.py --crossing-only
\end{quote}
using Python's standard library and a C++17 compiler. The default replay works
in a temporary directory and compares an immutable report. The universal
claims are the displayed analytic arguments, not extrapolations of the tests.

